\begin{theorem}[boundary 的双层遗忘与 \texorpdfstring{$56/97155$}{56/97155} 信息落差]\label{thm:conclusion-window6-boundary-double-forgetfulness-information-gap}\leanverified{paper\_conclusion\_window6\_boundary\_double\_forgetfulness\_information\_gap}
设
\[
\mathcal A_6:=\CC[\mathcal R_6^{\mathrm{bin}}],
\qquad
I_{\partial}\subset \mathcal A_6,
\qquad
Z_{\partial}\subset Z(G_6^{\mathrm{bin}})\cong (\ZZ/2\ZZ)^8
\]
分别为 canonical boundary 的三块理想与三维中心子格。则存在两级严格遗忘：
\begin{enumerate}
\normalfont
  \item 在抽象半单代数层面，
  \[
  \bigl|\operatorname{Orb}_{\Aut_{\CC\text{-alg}}(\mathcal A_6)}(I_{\partial})\bigr|
  =
  \binom{8}{3}
  =
  56;
  \]
  \item 在抽象群中心层面，
  \[
  \bigl|\operatorname{Orb}_{\Aut(G_6^{\mathrm{bin}})}(Z_{\partial})\bigr|
  =
  \binom{8}{3}_2
  =
  97155.
  \]
\end{enumerate}
故群层的歧义规模相较于代数层额外放大为
\[
\frac{97155}{56}\approx 1734.910714,
\qquad
\log_2\frac{97155}{56}\approx 10.7606457023\ \text{bits}.
\]
因此，window-\(6\) 的 boundary 并不是一个在抽象层面自然可见的对象；即使已知其在半单代数层对应 dyadic 三块，进一步恢复 canonical 中心三轴仍需额外的组合外置信息。
\end{theorem}
\begin{proof}
第一条正是定理 \ref{thm:conclusion-window6-semisimple-boundary-forgetfulness} 的内容：全部与 \(M_2(\CC)^{\oplus 3}\) 同构的双边理想在
\[
\Aut_{\CC\text{-alg}}(\mathcal A_6)
\]
作用下构成单一轨道，其大小恰为 \(\binom{8}{3}=56\)。

第二条则由定理 \ref{thm:conclusion-window6-boundary-center-seventeen-bit-barrier} 给出：
\[
\bigl|\operatorname{Orb}_{\Aut(G_6^{\mathrm{bin}})}(Z_{\partial})\bigr|
=
\binom{8}{3}_2
=
97155.
\]
两者相除并取二进对数，即得所述额外信息落差。证毕。
\end{proof}

\begin{corollary}[boundary parity 的三通道律与十五比特外置信息壁垒]\label{cor:conclusion-window6-boundary-threechannel-fifteenbit-externality}
\leanverified{paper\_conclusion\_window6\_boundary\_threechannel\_fifteenbit\_externality}
设 canonical boundary parity 的三类连续/几何通道秩分别为
\[
r_{\mathrm{rat}}=0,\qquad
r_{\mathrm{geo}}=1,\qquad
r_{\mathrm{faith}}=3.
\]
则 faithful boundary parity 的连续实现同时必须满足：
\begin{enumerate}
\normalfont
  \item 任一忠实 torus holonomy 实现至少需要三圆，即最小 torus 秩满足
  \[
  \operatorname{rank}T\ge 3;
  \]
  \item 即使 visible 根云 pinning 已把残余歧义压到 Klein 四量级，抽象中心恢复 canonical \(Z_{\partial}\) 仍至少需要
  \[
  17-2=15
  \]
  比特外置信息。
\end{enumerate}
因此，window-\(6\) 的 canonical boundary cube 不可能仅由 visible root cloud 的连续极限恢复；它在 faithful 通道中本质上是一份三维 torus 级别的离散外置账本。
\end{corollary}
\begin{proof}
由定理 \ref{thm:conclusion-window6-boundary-parity-zero-one-three-law}，
\[
r_{\mathrm{rat}}=0,\qquad r_{\mathrm{geo}}=1,\qquad r_{\mathrm{faith}}=3.
\]
其中 \(r_{\mathrm{faith}}=3\) 正说明：若要求忠实地以紧 torus holonomy 实现全部 boundary parity，则最小 torus 秩恰为 \(3\)，故必有 \(\operatorname{rank}T\ge 3\)。

另一方面，定理 \ref{thm:conclusion-window6-boundary-center-seventeen-bit-barrier} 已给出：从抽象群层精确恢复 canonical \(Z_{\partial}\) 至少需要 \(17\) 个二值选择，而 visible 根云 pinning 只剩 Klein 四量级的两位歧义。故额外所需外置信息至少为
\[
17-2=15
\]
比特。证毕。
\end{proof}

\begin{theorem}[动力学中心化子的 torus 塌缩]\label{thm:conclusion-window6-dynamical-centralizer-torus-collapse}\leanverified{paper\_conclusion\_window6\_dynamical\_centralizer\_torus\_collapse}
设 \(H\) 为作用于
\[
\mathcal H_6=\ell^2(X_6,\pi_6)
\]
的连通紧 Lie 群，并满足其表示 \(\rho(H)\) 与 pushforward 动力学算子 \(\mathsf T_6\) 对易。则经适当酉共轭后，
\[
\rho(H)\subset U(1)^{21},
\]
从而
\[
\Lie(\rho(H))
\]
必为阿贝尔 Lie 代数。特别地，任何 \(B_3/C_3\)-型非阿贝尔几何机制都不可能作为动力学中心化子的一部分存活。
\end{theorem}
\begin{proof}
这正是推论 \ref{cor:conclusion-window6-torus-phase-only-common-shadow} 的内容。该推论由推论 \ref{cor:conclusion-window6-dynamical-centralizer-toralization} 与定理 \ref{thm:conclusion-window6-dual-lie-realizations-incompatible} 联合推出：一切与 \(\mathsf T_6\) 对易的连通紧 Lie 群都可被酉共轭入 \(U(1)^{21}\)，而同载体上的几何 Lie 结构与动力学 Lie 包络又彼此不相容，故共同连通影像只能退化为环面相位。证毕。
\end{proof}

\begin{theorem}[最小壳层上 boundary/internal 的 Walsh 能量 \texorpdfstring{$25{:}9$}{25:9} 比]\label{thm:conclusion-window6-minshell-boundary-internal-walsh-energy-ratio}
\leanverified{paper\_conclusion\_window6\_minshell\_boundary\_internal\_walsh\_energy\_ratio}
在最小二重点壳
\[
Y_2\cong (\ZZ/2\ZZ)^3
\]
上，定义测度
\[
\mu_{\partial}:=\frac13\,\mathbf 1_{X_6^{\mathrm{bdry}}},
\qquad
\mu_{\mathrm{int}}:=\frac15\,\mathbf 1_{X_{6,\mathrm{int}}^{(2)}},
\qquad
u\equiv 2^{-3}.
\]
则碰撞概率满足
\[
\operatorname{Col}(\mu_{\partial})=\frac13,
\qquad
\operatorname{Col}(\mu_{\mathrm{int}})=\frac15.
\]
由此
\[
\chi^2(\mu_{\partial}\|u)=\frac53,
\qquad
\chi^2(\mu_{\mathrm{int}}\|u)=\frac35.
\]
等价地，两者的非平凡 Walsh 总能量之比精确为
\[
\frac{\sum_{a\neq 0}\widehat{\mu_{\partial}}(a)^2}
{\sum_{a\neq 0}\widehat{\mu_{\mathrm{int}}}(a)^2}
=
\frac{5/3}{3/5}
=
\frac{25}{9}.
\]
因此，在最小二重点壳层上，boundary 三点并非 internal 五点的一个稀薄补集；它在 Walsh 频域中精确携带 \(25/9\) 倍的非平凡二阶能量。
\end{theorem}
\begin{proof}
由推论 \ref{cor:conclusion-window6-twofiber-boundary-internal-splitting}，
\[
Y_2=X_6^{\mathrm{bdry}}\sqcup X_{6,\mathrm{int}}^{(2)},
\qquad
|Y_2|=8,\qquad |X_6^{\mathrm{bdry}}|=3,\qquad |X_{6,\mathrm{int}}^{(2)}|=5.
\]
故 \(\mu_{\partial}\) 与 \(\mu_{\mathrm{int}}\) 分别是在三点集与五点集上的均匀分布，于是
\[
\operatorname{Col}(\mu_{\partial})=\frac13,
\qquad
\operatorname{Col}(\mu_{\mathrm{int}})=\frac15.
\]

另一方面，由命题 \ref{prop:conclusion-window6-boundary-plancherel-information-identity}，
\[
\chi^2(\mu\|u)=8\Bigl(\operatorname{Col}(\mu)-\frac18\Bigr)
=
\sum_{a\neq 0}\widehat\mu(a)^2.
\]
代入 \(\mu=\mu_{\partial}\) 与 \(\mu=\mu_{\mathrm{int}}\) 分别得到
\[
\chi^2(\mu_{\partial}\|u)
=
8\left(\frac13-\frac18\right)
=
\frac53,
\]
\[
\chi^2(\mu_{\mathrm{int}}\|u)
=
8\left(\frac15-\frac18\right)
=
\frac35.
\]
因此
\[
\frac{\sum_{a\neq 0}\widehat{\mu_{\partial}}(a)^2}
{\sum_{a\neq 0}\widehat{\mu_{\mathrm{int}}}(a)^2}
=
\frac{\chi^2(\mu_{\partial}\|u)}{\chi^2(\mu_{\mathrm{int}}\|u)}
=
\frac{5/3}{3/5}
=
\frac{25}{9}.
\]
证毕。
\end{proof}

\begin{conjecture}[Walsh--Zeckendorf 对偶的归一化唯一实现]\label{conj:conclusion-window6-normalized-walsh-zeckendorf-duality}
在刚性识别
\[
S_{6,2}\cong 13+\{0,1,\dots,7\},
\qquad
\widehat Z_{\partial}\cong (\ZZ/2\ZZ)^3
\]
下，若进一步固定 boundary 三轴与三条坐标角色的对应，则应存在唯一的归一化酉 Fourier 变换
\[
\mathcal W_6:\CC[S_{6,2}]\longrightarrow \CC[\widehat Z_{\partial}]
\]
满足：
\begin{enumerate}
\normalfont
  \item \(\CC[X_6^{\mathrm{bdry}}]\) 被送到三条坐标角色张成的子空间；
  \item \(\CC[X_{6,\mathrm{cyc},2}]\) 被送到五个非坐标 Walsh 模张成的子空间；
  \item 强局域几何可实现的对角子群
  \[
  \langle(1,1,1)\rangle
  \]
  在频域上恰对应 \((4+4)\) 的粗粒二分。
\end{enumerate}
若此猜想成立，则 window-\(6\) 的最小壳层 \((5+3)\) 裂分、boundary parity 立方体以及 Walsh 模态层并非三套平行字典，而是同一有限 Fourier 对偶的三种坐标化。
\end{conjecture}
\begin{proof}[证据]
定理 \ref{thm:conclusion-window6-residual-threeinterval-terminal-window} 已把
\[
S_{6,2}
\]
刚性识别为终端 Zeckendorf 窗 \(13+\{0,\dots,7\}\)；定理 \ref{thm:conclusion-window6-minimal-shell-five-three-cleavage} 给出最小壳的 \((5+3)\) 裂分；定理 \ref{thm:conclusion-window6-walsh-spin-duality} 则把最小壳的八点结构与 Walsh--Spin 词典严格对应；而推论 \ref{cor:conclusion-window6-geometric-uplift-four-four-collapse} 又说明强局域几何在频域上留下精确的 \((4+4)\) 粗粒二分。

因此，终端 Zeckendorf 窗、最小壳的 \((5+3)\) 裂分与 boundary parity 角色立方体之间，已经存在一组与有限 Fourier 对偶高度一致的结构线索。尚未闭合的部分仅是显式且归一化唯一的 \(\mathcal W_6\) 构造，故保留为猜想。证毕。
\end{proof}

\paragraph{统一读法}
这一组结论表明，window-\(6\) 的 boundary 结构并不是若干彼此松散的附属现象。更严格的情形是：boundary 指示函数先在 radial 壳层代数上分裂出精确的 \(11{:}19\) 正交比例，并由六原子极小交换放大刻画其横截刚性；同标签几何随后表现为无边、混奇与长程三性并存，并在超公平噪声 \(p_*>1/2\) 处达到唯一再入极小值；三域 Frobenius 乘积密度则把 mixed cumulant 的消失提升为整个类函数 Hilbert 空间的 Hoeffding--ANOVA 分层；而在 boundary parity 一侧，抽象半单代数层的 \(56\) 重遗忘与抽象中心层的 \(97155\) 重遗忘共同把 faithful 恢复压成至少 \(15\) 比特的外置信息壁垒，并把 commuting dynamics 的共同连通影像彻底压缩为 \(U(1)^{21}\) 内的 torus phase shadow。由此，boundary 横截刚性、同标签再入噪声、三域正交独立与 canonical boundary cube 的外置恢复，便在同一终端 window-\(6\) 对象上闭合为一条统一结构链。

\endinput
